\section{The Principle of Mathematical Induction}

Many statements in mathematics are claims about every natural number at once: a formula that should hold for all $n$, an inequality valid from some point onward, a property shared by an infinite family of objects. Mathematical induction is the standard tool for proving such statements. We begin by stating the principle itself.

\begin{definition}[Principle of mathematical induction]
Let $P(n)$ be a statement defined for every integer $n \geq n_0$. Suppose that
\begin{enumerate}
    \item $P(n_0)$ is true, and
    \item for every integer $k \geq n_0$, the truth of $P(k)$ implies the truth of $P(k+1)$.
\end{enumerate}
Then $P(n)$ is true for all integers $n \geq n_0$.
\end{definition}

\subsection{The base case and the inductive step}

A proof by induction has exactly two parts, mirroring the two conditions above.

The \emph{base case} verifies the statement for the smallest value $n_0$, usually $n_0 = 1$. This anchors the argument; without it, nothing is established.

The \emph{inductive step} assumes the statement for an arbitrary $n = k$ and uses that assumption to deduce the statement for $n = k+1$. The assumption $P(k)$ is called the \emph{inductive hypothesis}. It is crucial to understand that we do not assume what we want to prove for all $n$; we assume it only for a single, arbitrary $k$, and then show the truth propagates one step forward.

\subsection{Intuition behind induction}

The usual picture is a line of dominoes. The base case is knocking over the first domino. The inductive step is the guarantee that each domino, when it falls, knocks over the next. Together these ensure that every domino eventually falls, no matter how far down the line it sits.

A complementary way to see why the principle is valid: suppose, for contradiction, that $P(n)$ fails for some $n$. Then there is a smallest value $m$ for which $P(m)$ is false. By the base case $m \neq n_0$, so $m - 1 \geq n_0$ and $P(m-1)$ is true (since $m$ was smallest). But the inductive step then forces $P(m)$ to be true, a contradiction. This is the well-ordering of the natural numbers doing the work behind the scenes.

\subsection{A first proof by induction}

We now prove the classic formula for the sum of the first $n$ positive integers.

\begin{theorem}\label{thm:sum}
For every integer $n \geq 1$,
\[
1 + 2 + \cdots + n = \frac{n(n+1)}{2}.
\]
\end{theorem}

\begin{proof}
Let $P(n)$ be the statement that $1 + 2 + \cdots + n = \dfrac{n(n+1)}{2}$.

\emph{Base case.} For $n = 1$, the left-hand side is $1$ and the right-hand side is $\dfrac{1 \cdot 2}{2} = 1$. The two agree, so $P(1)$ holds.

\emph{Inductive step.} Assume $P(k)$ holds for some integer $k \geq 1$; that is,
\[
1 + 2 + \cdots + k = \frac{k(k+1)}{2}.
\]
We must show $P(k+1)$, namely that $1 + 2 + \cdots + (k+1) = \dfrac{(k+1)(k+2)}{2}$. Adding $k+1$ to both sides of the inductive hypothesis gives
\[
1 + 2 + \cdots + k + (k+1) = \frac{k(k+1)}{2} + (k+1).
\]
Combining the terms on the right over a common denominator,
\[
\frac{k(k+1)}{2} + (k+1) = \frac{k(k+1) + 2(k+1)}{2} = \frac{(k+1)(k+2)}{2}.
\]
This is exactly $P(k+1)$.

By the principle of mathematical induction, $P(n)$ holds for all integers $n \geq 1$.
\end{proof}

\subsection{A second example}

To show the method is not limited to summation formulas, we prove a divisibility statement.

\begin{theorem}\label{thm:divisible}
For every integer $n \geq 1$, the number $n^3 - n$ is divisible by $6$.
\end{theorem}

\begin{proof}
Let $P(n)$ be the statement that $6 \mid (n^3 - n)$.

\emph{Base case.} For $n = 1$, we have $1^3 - 1 = 0$, and $6 \mid 0$, so $P(1)$ holds.

\emph{Inductive step.} Assume $P(k)$ holds for some integer $k \geq 1$, so that $k^3 - k$ is divisible by $6$. Consider
\[
(k+1)^3 - (k+1) = k^3 + 3k^2 + 3k + 1 - k - 1 = (k^3 - k) + 3k^2 + 3k.
\]
The first term $k^3 - k$ is divisible by $6$ by the inductive hypothesis. For the remaining part,
\[
3k^2 + 3k = 3k(k+1),
\]
and since $k$ and $k+1$ are consecutive integers one of them is even, so $k(k+1)$ is divisible by $2$ and $3k(k+1)$ is divisible by $6$. The sum of two multiples of $6$ is a multiple of $6$, so $(k+1)^3 - (k+1)$ is divisible by $6$. This is $P(k+1)$.

By the principle of mathematical induction, $6 \mid (n^3 - n)$ for all integers $n \geq 1$.
\end{proof}

\begin{remark}
Notice the shape shared by both proofs: verify a small case, assume the claim at $k$, and manipulate the expression at $k+1$ until the inductive hypothesis can be substituted in. Recognizing this pattern is most of the skill in writing inductive proofs.
\end{remark}
